\section{Persona Construction Method}

\subsection{Profile Database}

Each reviewer persona is constructed from a structured profile containing:

\begin{table}[h]
\centering
\small
\begin{tabular}{ll}
\toprule
Field & Example \\
\midrule
Name & Percy Liang \\
Affiliation & Stanford (HELM) \\
Category & ML Research / Learning \\
Expertise & HELM, benchmarks, foundations \\
Key Question & ``How was this evaluated?'' \\
Venue Alignment & ICML, NeurIPS, ICLR \\
Expertise Tags & \texttt{evaluation}, \texttt{benchmarks} \\
\bottomrule
\end{tabular}
\caption{Structured profile fields for persona construction.}
\end{table}

The database contains 45+ researchers across 10 categories, selected for coverage of major AI/ML research areas.

\subsection{Prompt Construction}

The review prompt for each persona includes:
\begin{enumerate}
    \item \textbf{Identity block}: ``You are [Name] from [Affiliation], an expert in [Expertise].''
    \item \textbf{Key question injection}: ``Your characteristic concern is: [Key Question]''
    \item \textbf{Venue context}: ``You are reviewing for [Venue], where papers are expected to [venue standards].''
    \item \textbf{Review structure}: Score rubric (1--4), required sections (major issues, minor issues, strengths, recommendations)
    \item \textbf{Paper content}: Full paper text or relevant sections
\end{enumerate}

\subsection{Panel Composition Rules}

Panels follow diversity constraints:
\begin{itemize}
    \item Minimum 5 reviewers per paper
    \item At least 1 industry practitioner and 1 academic
    \item At least 2 different expertise categories
    \item Maximum 2 reviewers from the same institution
    \item Venue-matched reviewers prioritized
\end{itemize}
